% 07_memory.tex -- DRAM controller, PHY, LPDDR chip, bank internals, timing, address map, bandwidth.
% Build: tectonic 07_memory.tex
\documentclass[tikz,border=10pt,12pt]{standalone}
% soc_style.tex -- shared look for the SoC figures (same palette as the slide decks and the
% core register placement figure). Every figure does % soc_style.tex -- shared look for the SoC figures (same palette as the slide decks and the
% core register placement figure). Every figure does % soc_style.tex -- shared look for the SoC figures (same palette as the slide decks and the
% core register placement figure). Every figure does \input{soc_style}.
\usepackage[T1]{fontenc}
\usepackage[scaled=0.95]{helvet}
\renewcommand{\familydefault}{\sfdefault}
\usepackage{booktabs}
\usepackage{amsmath}
\usepackage{amssymb}
\usetikzlibrary{arrows.meta,calc,positioning,fit,shapes.geometric,shapes.misc}

% ---------------- palette ----------------
\definecolor{pcb}{HTML}{0B2E33}
\definecolor{pcbmid}{HTML}{1F5C61}
\definecolor{tint}{HTML}{E6F0EF}
\definecolor{copper}{HTML}{B5651D}
\definecolor{coppert}{HTML}{F6E8DC}
\definecolor{ink}{HTML}{1F2A2E}
\definecolor{muted}{HTML}{5B6B70}
\definecolor{hardf}{HTML}{DCE6F2}
\definecolor{harde}{HTML}{2F4A7A}
\definecolor{good}{HTML}{2E7D32}
\definecolor{bad}{HTML}{B3261E}
% one colour per kind of block: border (d) / fill (f)
\definecolor{cpud}{HTML}{2B5C9E} \definecolor{cpuf}{HTML}{D6E4F5}
\definecolor{gpud}{HTML}{2E7D32} \definecolor{gpuf}{HTML}{D5EDD0}
\definecolor{npud}{HTML}{0B2E33} \definecolor{npuf}{HTML}{BFE0DB}
\definecolor{memd}{HTML}{5E3A8C} \definecolor{memf}{HTML}{E4DAF3}
\definecolor{busd}{HTML}{B5651D} \definecolor{busf}{HTML}{F6E8DC}
\definecolor{prd}{HTML}{546E7A}  \definecolor{prf}{HTML}{E3EAED}
\definecolor{secd}{HTML}{B3261E} \definecolor{secf}{HTML}{F7D9D6}
\definecolor{pwrd}{HTML}{8A6800} \definecolor{pwrf}{HTML}{FFF0B8}
\definecolor{medd}{HTML}{8E2463} \definecolor{medf}{HTML}{F5D6E6}
\definecolor{skd}{HTML}{2E7D32}  \definecolor{skf}{HTML}{BFE3B4}   % skew buffers, as in the register figure

\tikzset{
  blk/.style 2 args={draw=#1, fill=#2, rounded corners=4pt, thick, align=center,
                     font=\scriptsize, inner sep=3pt},
  cpu/.style={blk={cpud}{cpuf}},  gpu/.style={blk={gpud}{gpuf}},
  npu/.style={blk={npud}{npuf}},  mem/.style={blk={memd}{memf}},
  bus/.style={blk={busd}{busf}},  per/.style={blk={prd}{prf}},
  sec/.style={blk={secd}{secf}},  pwr/.style={blk={pwrd}{pwrf}},
  med/.style={blk={medd}{medf}},
  comb/.style ={draw=copper, fill=coppert, rounded corners=2pt, thick, align=center,
                font=\scriptsize, inner sep=2pt},                       % combinational logic
  hard/.style ={draw=harde, fill=hardf, thick, dashed, align=center,
                font=\scriptsize, inner sep=2pt},                       % hard macro / off-chip
  sram/.style ={draw=memd, fill=memf, thick, align=center, font=\scriptsize,
                inner sep=2pt},                                         % SRAM array, not flip-flops
  ff/.style   ={draw=#1, fill=#1!30, line width=0.35pt},                % one flip-flop square
  ffgone/.style={draw=muted, fill=white, densely dashed, line width=0.35pt},
  panel/.style={draw=muted, rounded corners=6pt, line width=0.9pt},
  ptitle/.style={anchor=north west, font=\normalsize\bfseries, text=ink, inner sep=0pt},
  note/.style ={font=\scriptsize, text=ink, align=left, inner sep=0pt},
  lab/.style  ={font=\tiny, text=muted, inner sep=1pt},
  axi/.style  ={{Stealth[length=2.6mm]}-{Stealth[length=2.6mm]}, line width=2.2pt, copper},
  axi1/.style ={-{Stealth[length=2.6mm]}, line width=2.2pt, copper},
  apb/.style  ={{Stealth[length=1.8mm]}-{Stealth[length=1.8mm]}, line width=1pt, copper},
  apb1/.style ={-{Stealth[length=1.8mm]}, line width=1pt, copper},
  dat/.style  ={-{Stealth[length=1.8mm]}, thick, pcbmid},
  ctl/.style  ={-{Stealth[length=1.8mm]}, densely dashed, thick, muted},
  irq/.style  ={-{Stealth[length=2mm]}, densely dotted, line width=1.3pt, bad},
  num/.style  ={circle, fill=copper, text=white, font=\scriptsize\bfseries, inner sep=1pt,
                minimum size=0.46cm},
  ours/.style ={draw=bad, line width=1.2pt, densely dashed, rounded corners=4pt},
}

% ---------------- helpers ----------------
\newcommand{\rn}[1]{\texttt{#1}}
% \ffrow{x}{y}{n}{cell}{style}: n flip-flops in a row, bottom-left corner at (x,y)
\newcommand{\ffrow}[5]{\foreach \q in {1,...,#3}{\draw[#5] ({#1+(\q-1)*#4},{#2}) rectangle ++(#4,#4);}}
% \ffcol{x}{y}{n}{cell}{style}: n flip-flops in a column, top-left corner at (x,y)
\newcommand{\ffcol}[5]{\foreach \q in {1,...,#3}{\draw[#5] ({#1},{#2-\q*#4}) rectangle ++(#4,#4);}}
% \ffgrid{x}{ytop}{n}{per row}{cell}{style}: n flip-flops, filled row by row from the top-left
\newcommand{\ffgrid}[6]{\pgfmathtruncatemacro{\gl}{#3-1}%
  \foreach \b in {0,...,\gl}{\pgfmathtruncatemacro{\gc}{mod(\b,#4)}\pgfmathtruncatemacro{\gr}{div(\b,#4)}%
    \draw[#6] ({#1+\gc*#5},{#2-(\gr+1)*#5}) rectangle ++(#5,#5);}}

% every figure: same page size, title, footer
\newcommand{\canvas}{\useasboundingbox (-0.3,-0.2) rectangle (38.6,30.4);}
\newcommand{\figtitle}[3]{%
  \node[num, minimum size=0.9cm, font=\Large\bfseries] at (0.45,29.55) {#1};
  \node[anchor=west, font=\LARGE\bfseries, text=pcb] at (1.15,29.55) {#2};
  \node[anchor=west, font=\small, text=ink] at (0,28.7) {#3};}
\newcommand{\figfoot}[1]{\node[anchor=west, font=\scriptsize, text=muted] at (0,0.25) {#1};}
% \panel{x0}{y0}{x1}{y1}{title}
\newcommand{\panel}[5]{\draw[panel] (#1,#2) rectangle (#3,#4);
  \node[ptitle] at ({#1+0.3},{#4-0.25}) {#5};}

% ---------------- timing diagrams: 1 cycle = \cw cm, low at y, high at y+0.45 ----------------
\newcommand{\cw}{1.0}
\newcommand{\wclk}[3]{% x0 y ncycles
  \foreach \k in {1,...,#3}{\draw[pcb, thick] ({#1+(\k-1)*\cw},#2) -- ++(0,0.45) -- ++(0.5*\cw,0)
     -- ++(0,-0.45) -- ++(0.5*\cw,0);}}
\newcommand{\wsig}[3]{% x0 y {level per cycle, e.g. 0,1,1,0}
  \gdef\wprev{x}%
  \foreach \v [count=\k from 0] in {#3}{%
    \draw[pcb, thick] ({#1+\k*\cw},{#2+0.45*\v}) -- ++(\cw,0);
    \if x\wprev\else\ifnum\wprev=\v\else\draw[pcb, thick] ({#1+\k*\cw},#2) -- ++(0,0.45);\fi\fi
    \global\let\wprev\v}}
\newcommand{\wbus}[5]{% x0 y c0 c1 text: one value held from cycle c0 to cycle c1
  \draw[pcbmid, thick, fill=white] ({#1+(#3)*\cw},{#2+0.225}) -- ({#1+(#3)*\cw+0.08},{#2+0.45})
    -- ({#1+(#4)*\cw-0.08},{#2+0.45}) -- ({#1+(#4)*\cw},{#2+0.225}) -- ({#1+(#4)*\cw-0.08},#2)
    -- ({#1+(#3)*\cw+0.08},#2) -- cycle;
  \node[font=\tiny\ttfamily] at ({#1+((#3)+(#4))*\cw/2},{#2+0.225}) {#5};}
\newcommand{\widle}[4]{% x0 y c0 c1: bus idle (no value)
  \draw[muted, thick] ({#1+(#3)*\cw},{#2+0.225}) -- ({#1+(#4)*\cw},{#2+0.225});}
\newcommand{\wname}[3]{\node[anchor=east, font=\scriptsize\ttfamily] at ({#1-0.15},{#2+0.225}) {#3};}
.
\usepackage[T1]{fontenc}
\usepackage[scaled=0.95]{helvet}
\renewcommand{\familydefault}{\sfdefault}
\usepackage{booktabs}
\usepackage{amsmath}
\usepackage{amssymb}
\usetikzlibrary{arrows.meta,calc,positioning,fit,shapes.geometric,shapes.misc}

% ---------------- palette ----------------
\definecolor{pcb}{HTML}{0B2E33}
\definecolor{pcbmid}{HTML}{1F5C61}
\definecolor{tint}{HTML}{E6F0EF}
\definecolor{copper}{HTML}{B5651D}
\definecolor{coppert}{HTML}{F6E8DC}
\definecolor{ink}{HTML}{1F2A2E}
\definecolor{muted}{HTML}{5B6B70}
\definecolor{hardf}{HTML}{DCE6F2}
\definecolor{harde}{HTML}{2F4A7A}
\definecolor{good}{HTML}{2E7D32}
\definecolor{bad}{HTML}{B3261E}
% one colour per kind of block: border (d) / fill (f)
\definecolor{cpud}{HTML}{2B5C9E} \definecolor{cpuf}{HTML}{D6E4F5}
\definecolor{gpud}{HTML}{2E7D32} \definecolor{gpuf}{HTML}{D5EDD0}
\definecolor{npud}{HTML}{0B2E33} \definecolor{npuf}{HTML}{BFE0DB}
\definecolor{memd}{HTML}{5E3A8C} \definecolor{memf}{HTML}{E4DAF3}
\definecolor{busd}{HTML}{B5651D} \definecolor{busf}{HTML}{F6E8DC}
\definecolor{prd}{HTML}{546E7A}  \definecolor{prf}{HTML}{E3EAED}
\definecolor{secd}{HTML}{B3261E} \definecolor{secf}{HTML}{F7D9D6}
\definecolor{pwrd}{HTML}{8A6800} \definecolor{pwrf}{HTML}{FFF0B8}
\definecolor{medd}{HTML}{8E2463} \definecolor{medf}{HTML}{F5D6E6}
\definecolor{skd}{HTML}{2E7D32}  \definecolor{skf}{HTML}{BFE3B4}   % skew buffers, as in the register figure

\tikzset{
  blk/.style 2 args={draw=#1, fill=#2, rounded corners=4pt, thick, align=center,
                     font=\scriptsize, inner sep=3pt},
  cpu/.style={blk={cpud}{cpuf}},  gpu/.style={blk={gpud}{gpuf}},
  npu/.style={blk={npud}{npuf}},  mem/.style={blk={memd}{memf}},
  bus/.style={blk={busd}{busf}},  per/.style={blk={prd}{prf}},
  sec/.style={blk={secd}{secf}},  pwr/.style={blk={pwrd}{pwrf}},
  med/.style={blk={medd}{medf}},
  comb/.style ={draw=copper, fill=coppert, rounded corners=2pt, thick, align=center,
                font=\scriptsize, inner sep=2pt},                       % combinational logic
  hard/.style ={draw=harde, fill=hardf, thick, dashed, align=center,
                font=\scriptsize, inner sep=2pt},                       % hard macro / off-chip
  sram/.style ={draw=memd, fill=memf, thick, align=center, font=\scriptsize,
                inner sep=2pt},                                         % SRAM array, not flip-flops
  ff/.style   ={draw=#1, fill=#1!30, line width=0.35pt},                % one flip-flop square
  ffgone/.style={draw=muted, fill=white, densely dashed, line width=0.35pt},
  panel/.style={draw=muted, rounded corners=6pt, line width=0.9pt},
  ptitle/.style={anchor=north west, font=\normalsize\bfseries, text=ink, inner sep=0pt},
  note/.style ={font=\scriptsize, text=ink, align=left, inner sep=0pt},
  lab/.style  ={font=\tiny, text=muted, inner sep=1pt},
  axi/.style  ={{Stealth[length=2.6mm]}-{Stealth[length=2.6mm]}, line width=2.2pt, copper},
  axi1/.style ={-{Stealth[length=2.6mm]}, line width=2.2pt, copper},
  apb/.style  ={{Stealth[length=1.8mm]}-{Stealth[length=1.8mm]}, line width=1pt, copper},
  apb1/.style ={-{Stealth[length=1.8mm]}, line width=1pt, copper},
  dat/.style  ={-{Stealth[length=1.8mm]}, thick, pcbmid},
  ctl/.style  ={-{Stealth[length=1.8mm]}, densely dashed, thick, muted},
  irq/.style  ={-{Stealth[length=2mm]}, densely dotted, line width=1.3pt, bad},
  num/.style  ={circle, fill=copper, text=white, font=\scriptsize\bfseries, inner sep=1pt,
                minimum size=0.46cm},
  ours/.style ={draw=bad, line width=1.2pt, densely dashed, rounded corners=4pt},
}

% ---------------- helpers ----------------
\newcommand{\rn}[1]{\texttt{#1}}
% \ffrow{x}{y}{n}{cell}{style}: n flip-flops in a row, bottom-left corner at (x,y)
\newcommand{\ffrow}[5]{\foreach \q in {1,...,#3}{\draw[#5] ({#1+(\q-1)*#4},{#2}) rectangle ++(#4,#4);}}
% \ffcol{x}{y}{n}{cell}{style}: n flip-flops in a column, top-left corner at (x,y)
\newcommand{\ffcol}[5]{\foreach \q in {1,...,#3}{\draw[#5] ({#1},{#2-\q*#4}) rectangle ++(#4,#4);}}
% \ffgrid{x}{ytop}{n}{per row}{cell}{style}: n flip-flops, filled row by row from the top-left
\newcommand{\ffgrid}[6]{\pgfmathtruncatemacro{\gl}{#3-1}%
  \foreach \b in {0,...,\gl}{\pgfmathtruncatemacro{\gc}{mod(\b,#4)}\pgfmathtruncatemacro{\gr}{div(\b,#4)}%
    \draw[#6] ({#1+\gc*#5},{#2-(\gr+1)*#5}) rectangle ++(#5,#5);}}

% every figure: same page size, title, footer
\newcommand{\canvas}{\useasboundingbox (-0.3,-0.2) rectangle (38.6,30.4);}
\newcommand{\figtitle}[3]{%
  \node[num, minimum size=0.9cm, font=\Large\bfseries] at (0.45,29.55) {#1};
  \node[anchor=west, font=\LARGE\bfseries, text=pcb] at (1.15,29.55) {#2};
  \node[anchor=west, font=\small, text=ink] at (0,28.7) {#3};}
\newcommand{\figfoot}[1]{\node[anchor=west, font=\scriptsize, text=muted] at (0,0.25) {#1};}
% \panel{x0}{y0}{x1}{y1}{title}
\newcommand{\panel}[5]{\draw[panel] (#1,#2) rectangle (#3,#4);
  \node[ptitle] at ({#1+0.3},{#4-0.25}) {#5};}

% ---------------- timing diagrams: 1 cycle = \cw cm, low at y, high at y+0.45 ----------------
\newcommand{\cw}{1.0}
\newcommand{\wclk}[3]{% x0 y ncycles
  \foreach \k in {1,...,#3}{\draw[pcb, thick] ({#1+(\k-1)*\cw},#2) -- ++(0,0.45) -- ++(0.5*\cw,0)
     -- ++(0,-0.45) -- ++(0.5*\cw,0);}}
\newcommand{\wsig}[3]{% x0 y {level per cycle, e.g. 0,1,1,0}
  \gdef\wprev{x}%
  \foreach \v [count=\k from 0] in {#3}{%
    \draw[pcb, thick] ({#1+\k*\cw},{#2+0.45*\v}) -- ++(\cw,0);
    \if x\wprev\else\ifnum\wprev=\v\else\draw[pcb, thick] ({#1+\k*\cw},#2) -- ++(0,0.45);\fi\fi
    \global\let\wprev\v}}
\newcommand{\wbus}[5]{% x0 y c0 c1 text: one value held from cycle c0 to cycle c1
  \draw[pcbmid, thick, fill=white] ({#1+(#3)*\cw},{#2+0.225}) -- ({#1+(#3)*\cw+0.08},{#2+0.45})
    -- ({#1+(#4)*\cw-0.08},{#2+0.45}) -- ({#1+(#4)*\cw},{#2+0.225}) -- ({#1+(#4)*\cw-0.08},#2)
    -- ({#1+(#3)*\cw+0.08},#2) -- cycle;
  \node[font=\tiny\ttfamily] at ({#1+((#3)+(#4))*\cw/2},{#2+0.225}) {#5};}
\newcommand{\widle}[4]{% x0 y c0 c1: bus idle (no value)
  \draw[muted, thick] ({#1+(#3)*\cw},{#2+0.225}) -- ({#1+(#4)*\cw},{#2+0.225});}
\newcommand{\wname}[3]{\node[anchor=east, font=\scriptsize\ttfamily] at ({#1-0.15},{#2+0.225}) {#3};}
.
\usepackage[T1]{fontenc}
\usepackage[scaled=0.95]{helvet}
\renewcommand{\familydefault}{\sfdefault}
\usepackage{booktabs}
\usepackage{amsmath}
\usepackage{amssymb}
\usetikzlibrary{arrows.meta,calc,positioning,fit,shapes.geometric,shapes.misc}

% ---------------- palette ----------------
\definecolor{pcb}{HTML}{0B2E33}
\definecolor{pcbmid}{HTML}{1F5C61}
\definecolor{tint}{HTML}{E6F0EF}
\definecolor{copper}{HTML}{B5651D}
\definecolor{coppert}{HTML}{F6E8DC}
\definecolor{ink}{HTML}{1F2A2E}
\definecolor{muted}{HTML}{5B6B70}
\definecolor{hardf}{HTML}{DCE6F2}
\definecolor{harde}{HTML}{2F4A7A}
\definecolor{good}{HTML}{2E7D32}
\definecolor{bad}{HTML}{B3261E}
% one colour per kind of block: border (d) / fill (f)
\definecolor{cpud}{HTML}{2B5C9E} \definecolor{cpuf}{HTML}{D6E4F5}
\definecolor{gpud}{HTML}{2E7D32} \definecolor{gpuf}{HTML}{D5EDD0}
\definecolor{npud}{HTML}{0B2E33} \definecolor{npuf}{HTML}{BFE0DB}
\definecolor{memd}{HTML}{5E3A8C} \definecolor{memf}{HTML}{E4DAF3}
\definecolor{busd}{HTML}{B5651D} \definecolor{busf}{HTML}{F6E8DC}
\definecolor{prd}{HTML}{546E7A}  \definecolor{prf}{HTML}{E3EAED}
\definecolor{secd}{HTML}{B3261E} \definecolor{secf}{HTML}{F7D9D6}
\definecolor{pwrd}{HTML}{8A6800} \definecolor{pwrf}{HTML}{FFF0B8}
\definecolor{medd}{HTML}{8E2463} \definecolor{medf}{HTML}{F5D6E6}
\definecolor{skd}{HTML}{2E7D32}  \definecolor{skf}{HTML}{BFE3B4}   % skew buffers, as in the register figure

\tikzset{
  blk/.style 2 args={draw=#1, fill=#2, rounded corners=4pt, thick, align=center,
                     font=\scriptsize, inner sep=3pt},
  cpu/.style={blk={cpud}{cpuf}},  gpu/.style={blk={gpud}{gpuf}},
  npu/.style={blk={npud}{npuf}},  mem/.style={blk={memd}{memf}},
  bus/.style={blk={busd}{busf}},  per/.style={blk={prd}{prf}},
  sec/.style={blk={secd}{secf}},  pwr/.style={blk={pwrd}{pwrf}},
  med/.style={blk={medd}{medf}},
  comb/.style ={draw=copper, fill=coppert, rounded corners=2pt, thick, align=center,
                font=\scriptsize, inner sep=2pt},                       % combinational logic
  hard/.style ={draw=harde, fill=hardf, thick, dashed, align=center,
                font=\scriptsize, inner sep=2pt},                       % hard macro / off-chip
  sram/.style ={draw=memd, fill=memf, thick, align=center, font=\scriptsize,
                inner sep=2pt},                                         % SRAM array, not flip-flops
  ff/.style   ={draw=#1, fill=#1!30, line width=0.35pt},                % one flip-flop square
  ffgone/.style={draw=muted, fill=white, densely dashed, line width=0.35pt},
  panel/.style={draw=muted, rounded corners=6pt, line width=0.9pt},
  ptitle/.style={anchor=north west, font=\normalsize\bfseries, text=ink, inner sep=0pt},
  note/.style ={font=\scriptsize, text=ink, align=left, inner sep=0pt},
  lab/.style  ={font=\tiny, text=muted, inner sep=1pt},
  axi/.style  ={{Stealth[length=2.6mm]}-{Stealth[length=2.6mm]}, line width=2.2pt, copper},
  axi1/.style ={-{Stealth[length=2.6mm]}, line width=2.2pt, copper},
  apb/.style  ={{Stealth[length=1.8mm]}-{Stealth[length=1.8mm]}, line width=1pt, copper},
  apb1/.style ={-{Stealth[length=1.8mm]}, line width=1pt, copper},
  dat/.style  ={-{Stealth[length=1.8mm]}, thick, pcbmid},
  ctl/.style  ={-{Stealth[length=1.8mm]}, densely dashed, thick, muted},
  irq/.style  ={-{Stealth[length=2mm]}, densely dotted, line width=1.3pt, bad},
  num/.style  ={circle, fill=copper, text=white, font=\scriptsize\bfseries, inner sep=1pt,
                minimum size=0.46cm},
  ours/.style ={draw=bad, line width=1.2pt, densely dashed, rounded corners=4pt},
}

% ---------------- helpers ----------------
\newcommand{\rn}[1]{\texttt{#1}}
% \ffrow{x}{y}{n}{cell}{style}: n flip-flops in a row, bottom-left corner at (x,y)
\newcommand{\ffrow}[5]{\foreach \q in {1,...,#3}{\draw[#5] ({#1+(\q-1)*#4},{#2}) rectangle ++(#4,#4);}}
% \ffcol{x}{y}{n}{cell}{style}: n flip-flops in a column, top-left corner at (x,y)
\newcommand{\ffcol}[5]{\foreach \q in {1,...,#3}{\draw[#5] ({#1},{#2-\q*#4}) rectangle ++(#4,#4);}}
% \ffgrid{x}{ytop}{n}{per row}{cell}{style}: n flip-flops, filled row by row from the top-left
\newcommand{\ffgrid}[6]{\pgfmathtruncatemacro{\gl}{#3-1}%
  \foreach \b in {0,...,\gl}{\pgfmathtruncatemacro{\gc}{mod(\b,#4)}\pgfmathtruncatemacro{\gr}{div(\b,#4)}%
    \draw[#6] ({#1+\gc*#5},{#2-(\gr+1)*#5}) rectangle ++(#5,#5);}}

% every figure: same page size, title, footer
\newcommand{\canvas}{\useasboundingbox (-0.3,-0.2) rectangle (38.6,30.4);}
\newcommand{\figtitle}[3]{%
  \node[num, minimum size=0.9cm, font=\Large\bfseries] at (0.45,29.55) {#1};
  \node[anchor=west, font=\LARGE\bfseries, text=pcb] at (1.15,29.55) {#2};
  \node[anchor=west, font=\small, text=ink] at (0,28.7) {#3};}
\newcommand{\figfoot}[1]{\node[anchor=west, font=\scriptsize, text=muted] at (0,0.25) {#1};}
% \panel{x0}{y0}{x1}{y1}{title}
\newcommand{\panel}[5]{\draw[panel] (#1,#2) rectangle (#3,#4);
  \node[ptitle] at ({#1+0.3},{#4-0.25}) {#5};}

% ---------------- timing diagrams: 1 cycle = \cw cm, low at y, high at y+0.45 ----------------
\newcommand{\cw}{1.0}
\newcommand{\wclk}[3]{% x0 y ncycles
  \foreach \k in {1,...,#3}{\draw[pcb, thick] ({#1+(\k-1)*\cw},#2) -- ++(0,0.45) -- ++(0.5*\cw,0)
     -- ++(0,-0.45) -- ++(0.5*\cw,0);}}
\newcommand{\wsig}[3]{% x0 y {level per cycle, e.g. 0,1,1,0}
  \gdef\wprev{x}%
  \foreach \v [count=\k from 0] in {#3}{%
    \draw[pcb, thick] ({#1+\k*\cw},{#2+0.45*\v}) -- ++(\cw,0);
    \if x\wprev\else\ifnum\wprev=\v\else\draw[pcb, thick] ({#1+\k*\cw},#2) -- ++(0,0.45);\fi\fi
    \global\let\wprev\v}}
\newcommand{\wbus}[5]{% x0 y c0 c1 text: one value held from cycle c0 to cycle c1
  \draw[pcbmid, thick, fill=white] ({#1+(#3)*\cw},{#2+0.225}) -- ({#1+(#3)*\cw+0.08},{#2+0.45})
    -- ({#1+(#4)*\cw-0.08},{#2+0.45}) -- ({#1+(#4)*\cw},{#2+0.225}) -- ({#1+(#4)*\cw-0.08},#2)
    -- ({#1+(#3)*\cw+0.08},#2) -- cycle;
  \node[font=\tiny\ttfamily] at ({#1+((#3)+(#4))*\cw/2},{#2+0.225}) {#5};}
\newcommand{\widle}[4]{% x0 y c0 c1: bus idle (no value)
  \draw[muted, thick] ({#1+(#3)*\cw},{#2+0.225}) -- ({#1+(#4)*\cw},{#2+0.225});}
\newcommand{\wname}[3]{\node[anchor=east, font=\scriptsize\ttfamily] at ({#1-0.15},{#2+0.225}) {#3};}

\begin{document}
\begin{tikzpicture}[line join=round]
\canvas
\figtitle{7}{Memory system: DRAM controller, PHY and the LPDDR chip}
  {DRAM is big and cheap but slow: every access must first open a row. The controller reorders requests to open as few
   rows as possible; the PHY moves the bits at gigabit speed per pin.}

% =====================================================================
% (A) the path
% =====================================================================
\panel{0}{19.6}{38.3}{28.1}{From the NoC to a DRAM cell}
% NoC ports
\node[bus, minimum width=2.1cm, minimum height=6.4cm, text width=1.9cm, anchor=north west, font=\tiny] at (0.3,27.2)
  {\textbf{4 AXI ports}\\from the NoC\\[3pt]CPU\\GPU\\NPU\\display};
% front-end
\draw[mem] (2.7,20.0) rectangle (12.4,27.3);
\node[anchor=north west, font=\scriptsize\bfseries, text=memd] at (2.85,27.15) {controller front-end};
\node[comb, minimum width=2.7cm, minimum height=2.3cm, text width=2.5cm, anchor=north west, font=\tiny] at (2.9,26.6)
  {\textbf{port arbiter + QoS}\\display is real-time: it goes first, or the screen flickers};
\node[anchor=north west, font=\tiny\bfseries] at (5.85,26.6) {read queue, 32 entries};
\ffgrid{5.9}{26.25}{128}{16}{0.15}{ff=memd}
\node[anchor=north west, font=\tiny\bfseries] at (5.85,24.85) {write queue, 32 entries};
\ffgrid{5.9}{24.5}{128}{16}{0.15}{ff=memd}
\node[lab, anchor=north west, text width=2.6cm] at (5.85,23.1) {(one square = 1 entry field) address, ID, age: flip-flops};
\node[sram, minimum width=2.6cm, minimum height=0.7cm, anchor=north west, font=\tiny] at (5.85,22.2) {write data buffer (SRAM)};
\node[comb, minimum width=3.3cm, minimum height=2.3cm, text width=3.1cm, anchor=north west, font=\tiny] at (8.95,26.6)
  {\textbf{address mapper}\\physical address $\rightarrow$ channel, bank, row, column (panel D)};
\node[note, font=\tiny, anchor=north west, text width=3.3cm] at (8.95,24.0) {%
  Queues let the controller look at many requests and choose a good order instead of serving them one by one.};
% back-end
\draw[mem] (12.7,20.0) rectangle (24.6,27.3);
\node[anchor=north west, font=\scriptsize\bfseries, text=memd] at (12.85,27.15) {controller back-end};
\node[comb, minimum width=3.0cm, minimum height=2.3cm, text width=2.8cm, anchor=north west, font=\tiny] at (12.9,26.6)
  {\textbf{scheduler} (FR-FCFS)\\1. requests to an open row first\\2. then the oldest\\+ an age limit so nobody starves};
\node[anchor=north west, font=\tiny\bfseries] at (16.1,26.6) {16 bank state machines};
\foreach \b in {0,...,15}{
  \pgfmathtruncatemacro{\bc}{mod(\b,4)}\pgfmathtruncatemacro{\br}{div(\b,4)}
  \node[mem, fill=white, minimum width=1.0cm, minimum height=0.55cm, inner sep=0pt, font=\tiny, anchor=north west]
    at ({16.1+\bc*1.1},{26.2-\br*0.65}) {bank \b};}
\node[lab, anchor=north west, text width=4.3cm] at (16.05,23.55) {each: state (idle / opening / open / closing) + open-row number register};
\node[anchor=north west, font=\tiny\bfseries] at (20.65,26.6) {timing counters};
\foreach \t [count=\k from 0] in {tRCD, tRP, tRAS, tWR, tFAW}{
  \node[font=\tiny\ttfamily, anchor=west] at (20.65,{26.05-\k*0.42}) {\t};
  \ffrow{21.55}{25.95-\k*0.42}{8}{0.18}{ff=memd}}
\node[lab, anchor=north west, text width=3.8cm] at (20.6,23.85) {down-counters: a command may go only when its counter is 0};
\node[comb, minimum width=3.7cm, minimum height=1.4cm, text width=3.5cm, anchor=north west, font=\tiny] at (20.7,22.95)
  {\textbf{refresh timer}: every $\approx$3.9 $\mu$s send REF, so every row is refreshed within 32 ms};
\node[comb, minimum width=7.6cm, minimum height=0.8cm, anchor=north west, font=\tiny] at (12.9,21.3)
  {commands out: ACT, RD, WR, PRE, REF $\rightarrow$ DFI interface (standard controller $\leftrightarrow$ PHY)};
% PHY
\draw[hard] (24.9,20.0) rectangle (31.9,27.3);
\node[anchor=north west, font=\scriptsize\bfseries, text=harde] at (25.05,27.15) {PHY (mixed analog / digital)};
\foreach \nm/\tx [count=\k from 0] in {
    serializer/{8 bits per pin per controller clock (800 MHz) $\rightarrow$ 6.4 Gb/s on the pin},
    delay lines + training/{shift each data bit so it lines up with the strobe; tuned at boot},
    {I/O drivers}/{push the signal down the board trace; on-die termination stops reflections}}{
  \node[draw=harde, fill=white, rounded corners=3pt, minimum width=6.6cm, minimum height=1.55cm, text width=6.4cm, align=center,
        anchor=north west, font=\tiny] at (25.1,{26.6-\k*1.95}) {\textbf{\nm}\\\tx};}
% DRAM chip
\draw[hard] (32.3,20.0) rectangle (38.0,27.3);
\node[anchor=north west, font=\scriptsize\bfseries, text=harde] at (32.45,27.15) {LPDDR5 chip (off-chip)};
\foreach \nm/\tx [count=\k from 0] in {
    command decoder/{reads ACT / RD / PRE},
    16 banks/{4 bank groups $\times$ 4},
    {I/O}/{16 data pins + strobes per channel}}{
  \node[draw=harde, fill=white, rounded corners=3pt, minimum width=5.3cm, minimum height=1.55cm, text width=5.1cm, align=center,
        anchor=north west, font=\tiny] at (32.45,{26.6-\k*1.95}) {\textbf{\nm}\\\tx};}
\draw[axi1] (2.4,24.0) -- (2.7,24.0);
\draw[axi1] (12.4,24.0) -- (12.7,24.0);
\draw[axi1] (24.6,24.0) -- (24.9,24.0);
\draw[axi] (31.9,24.0) -- (32.3,24.0);

% =====================================================================
% (B) inside a DRAM bank
% =====================================================================
\panel{0}{0.7}{13.5}{19.3}{Inside one DRAM bank}
% one cell
\node[anchor=north west, font=\scriptsize\bfseries] at (0.3,18.6) {one cell = 1 transistor + 1 capacitor};
\draw[pcb, thick] (1.6,18.0) -- (1.6,14.6);
\node[lab, text=ink, anchor=north] at (1.6,14.55) {bitline (column)};
\draw[pcb, thick] (0.6,17.4) -- (4.4,17.4);
\node[lab, text=ink, anchor=south] at (3.7,17.42) {wordline (row)};
\draw[pcb, thick] (1.6,16.4) -- (2.3,16.4);
\draw[pcb, very thick] (2.3,16.15) -- (2.3,16.65);
\draw[pcb, very thick] (2.45,16.15) -- (2.45,16.65);
\draw[pcb, thick] (2.45,16.4) -- (3.2,16.4) -- (3.2,15.95);
\draw[pcb, very thick] (2.9,15.95) -- (3.5,15.95);
\draw[pcb, very thick] (2.9,15.75) -- (3.5,15.75);
\draw[pcb, thick] (3.2,15.75) -- (3.2,15.4);
\draw[pcb, thick] (2.95,15.4) -- (3.45,15.4) (3.05,15.3) -- (3.35,15.3) (3.15,15.2) -- (3.25,15.2);
\draw[pcb, thick] (2.3,16.65) -- (2.3,17.4);
\node[lab, text=ink, anchor=west] at (3.6,15.95) {capacitor};
\node[lab, text=ink, anchor=west] at (3.6,15.65) {charge = 1 bit};
\node[note, font=\scriptsize, anchor=north west, text width=4.6cm] at (0.3,14.15) {%
  The charge leaks away in milliseconds, so every row must be refreshed. Reading also drains it, so the row is written back.
  SRAM uses 6 transistors per bit and needs no refresh: faster but much bigger.};
% bank array
\node[anchor=north west, font=\scriptsize\bfseries] at (5.3,18.6) {bank = rows $\times$ columns of cells};
\node[comb, minimum width=0.6cm, minimum height=5.0cm, anchor=north west, font=\tiny, text width=0.5cm, align=center] at (5.3,18.15)
  {r\\o\\w\\[2pt]d\\e\\c};
\foreach \r in {0,...,9}{\foreach \c in {0,...,13}{
  \pgfmathsetmacro{\fillc}{\r==3 ? 1 : 0}
  \ifdim\fillc pt>0.5pt \draw[medd, fill=medf] ({6.05+\c*0.5},{17.7-\r*0.5}) rectangle ++(0.42,0.42);
  \else \draw[memd, fill=memf] ({6.05+\c*0.5},{17.7-\r*0.5}) rectangle ++(0.42,0.42); \fi}}
\node[lab, text=medd, font=\tiny\bfseries, anchor=west] at (13.0,16.4) {$\leftarrow$};
\node[comb, minimum width=7.0cm, minimum height=0.6cm, anchor=north west, font=\tiny] at (6.0,12.95)
  {\textbf{sense amplifiers = row buffer}: holds one whole open row (2 KB)};
\node[comb, minimum width=7.0cm, minimum height=0.6cm, anchor=north west, font=\tiny] at (6.0,12.2)
  {\textbf{column mux}: picks 32 bytes = one burst};
\draw[dat] (9.5,11.6) -- (9.5,11.1);
\node[lab, text=ink, anchor=north] at (9.5,11.05) {to the 16 data pins, 16 beats};
\node[note, anchor=north west, text width=12.8cm] at (0.3,10.3) {%
  \textbf{1 ACT row}: the row decoder raises one wordline; the whole row flows into the sense amplifiers (time tRCD).\\[3pt]
  \textbf{2 RD column}: the column mux picks 32 bytes from the open row; they reach the pins after the read latency,
  as 16 beats of 16 bits.\\[3pt]
  \textbf{3} More reads to the same open row are quick: a \textbf{row hit}.\\[3pt]
  \textbf{4 PRE}: close the row (write it back) before another row of this bank can open (time tRP).\\[3pt]
  \textbf{5 REF}: refresh rows regularly, or the data is lost.};

% =====================================================================
% (C) timing
% =====================================================================
\panel{13.8}{10.2}{38.3}{19.3}{Command timing on one bank (not to scale)}
\renewcommand{\cw}{0.95}
\wname{16.3}{17.6}{CK}  \wclk{16.3}{17.6}{22}
\wname{16.3}{16.75}{CMD}
\wbus{16.3}{16.75}{1}{2}{ACT} \widle{16.3}{16.75}{2}{6} \wbus{16.3}{16.75}{6}{7}{RD}
\widle{16.3}{16.75}{7}{16} \wbus{16.3}{16.75}{16}{17}{PRE} \widle{16.3}{16.75}{17}{20} \wbus{16.3}{16.75}{20}{21}{ACT} \widle{16.3}{16.75}{0}{1}
\wname{16.3}{15.9}{DQ}
\widle{16.3}{15.9}{0}{11} \wbus{16.3}{15.9}{11}{15}{16 data beats} \widle{16.3}{15.9}{15}{22}
\foreach \a/\b/\t in {1/6/tRCD, 6/11/read latency, 16/20/tRP}{
  \draw[{Stealth[length=1.5mm]}-{Stealth[length=1.5mm]}, bad, thick] ({16.3+\a*\cw},15.35) -- ({16.3+\b*\cw},15.35);
  \node[font=\tiny\bfseries, text=bad, anchor=north] at ({16.3+(\a+\b)*\cw/2},15.3) {\t};}
\node[lab, text=ink, anchor=south] at ({16.3+1.5*\cw},17.25) {open row 0x1A2};
\node[lab, text=ink, anchor=south] at ({16.3+6.5*\cw},17.25) {column 0x40};
\node[lab, text=ink, anchor=south] at ({16.3+20.5*\cw},17.25) {next row};
\node[anchor=north west, font=\scriptsize, inner sep=0pt] at (14.1,14.4) {%
  \renewcommand{\arraystretch}{1.15}%
  \begin{tabular}{@{}l l l@{}}
    \toprule
    \textbf{Case} & \textbf{What must happen} & \textbf{DRAM time (typical)}\\
    \midrule
    row hit      & the row is already open: RD only               & read latency $\approx$15--20 ns\\
    row closed   & ACT, then RD                                   & tRCD + latency $\approx$30--40 ns\\
    row conflict & PRE the old row, ACT the new one, then RD      & tRP + tRCD + latency $\approx$45--60 ns\\
    \bottomrule
  \end{tabular}};
\node[note, anchor=north west, text width=7.9cm] at (30.2,14.4) {%
  Add the controller queues, PHY and NoC and a CPU sees $\approx$100 ns for a DRAM read.\\[3pt]
  That is why the scheduler prefers row hits, and why reading memory in order (like the NPU's DMA bursts)
  is much faster than jumping around.};

% =====================================================================
% (D) address mapping + bandwidth
% =====================================================================
\panel{13.8}{0.7}{38.3}{9.9}{Address mapping and bandwidth}
\node[anchor=west, font=\scriptsize\bfseries] at (14.1,8.8) {example physical address split (4 channels, 16 banks, 2 KB rows, 32-byte bursts):};
\foreach \nm/\bits/\w/\col [count=\k from 0] in {row/{[33:17]}/6.0/cpuf, bank/{[16:13]}/2.2/pwrf, column/{[12:7]}/2.6/memf, channel/{[6:5]}/1.8/gpuf, byte/{[4:0]}/2.2/medf}{
  \ifnum\k=0 \xdef\ax{14.1}\fi
  \node[draw=pcb, fill=\col, minimum width=\w cm, minimum height=0.7cm, anchor=north west, inner sep=0pt, font=\tiny, align=center]
    at (\ax,8.35) {\textbf{\nm} \bits};
  \pgfmathsetmacro{\nx}{\ax+\w}\xdef\ax{\nx}}
\node[note, font=\tiny, anchor=north west, text width=9.0cm] at (29.2,8.4) {%
  Channel bits sit low, so neighbouring 32-byte pieces go to different channels and all 4 work at once.
  Real controllers also XOR some bits to spread traffic.};
\node[note, anchor=north west, text width=11.8cm] at (14.1,7.25) {%
  \textbf{Bandwidth} = transfers per second $\times$ bus width\\[2pt]
  LPDDR5-6400: 6,400 M/s $\times$ 16 bit = 102.4 Gb/s = \textbf{12.8 GB/s per channel}\\
  4 channels $\Rightarrow$ \textbf{51.2 GB/s} peak (real traffic gets less).\\[4pt]
  DE10-Lite: 64 MB SDRAM, 16 bit, single data rate. At 100 MHz: 16 bit $\times$ 100 M = \textbf{200 MB/s}.};
\node[note, anchor=north west, text width=11.8cm] at (26.2,7.25) {%
  \textbf{Why an NPU needs on-chip SRAM}\\[2pt]
  32$\times$32 MACs at 1 GHz, each needing 2 new bytes: 2,048 B per clock = \textbf{2 TB/s}.\\
  The systolic array reuses each value along a whole row or column: only 64 B per clock = \textbf{64 GB/s}.\\
  Still more than 51 GB/s $\Rightarrow$ tiles stay in SRAM and are reused many times (Figure 5).
  Your 4$\times$4 array uses every A and B value in 4 PEs.};

\figfoot{LPDDR5 numbers are typical; exact timings depend on the speed grade. The address split is an example; each controller chooses its own.}
\end{tikzpicture}
\end{document}
